\chapter{Tribal Communities in India}

\section{The Definitional Problem}

\subsection{Four Themes of the Tribal Unit}

Tribes constitute a vast area of interest --- almost a full discipline in themselves --- but in the syllabus they form one small unit from which one or two questions are asked.

\begin{KeyConcept}
\begin{itemize}
\item The \textbf{four major themes} under Tribes:
\item 1. The \textbf{definitional problem} (what is a tribe, and by what criteria are people labelled tribes?).
\item 2. \textbf{Geographical spread} of tribes.
\item 3. \textbf{Colonial policies and the tribes} (and their impact).
\item 4. The issue of \textbf{integration and autonomy}.
\item These themes are not isolated --- they are all integrated, and the root of everything lies in the \textbf{colonial history of India}.
\end{itemize}
\end{KeyConcept}

\subsection{Tribe as an Externally Imposed Category}

Do tribes see themselves as 'tribes' the way we identify them? No --- they do not identify themselves as such.

\begin{ExampleBox}
\begin{itemize}
\item \textbf{S. C. Roy} (India's first ethnographer) studied the \textbf{Mundas} and found that they call themselves \textbf{Horoko} --- a term in the Mundari language meaning \textbf{'the original creation of the creator'} --- not a subdivision of the population, a different race, or something lesser, but the original human beings.
\item When Roy interacted with the \textbf{Khonds} (a tribe), one of them asked, 'Are there Khonds in your country as well?' --- showing that the groups we treat as distinct view themselves completely differently from how outsiders see them.
\end{itemize}
\end{ExampleBox}

\begin{KeyConcept}
\begin{itemize}
\item The \textbf{tribal identity} by which these groups are now known is \textbf{not something they imposed on themselves or created} --- it is an \textbf{external category imposed upon them}.
\item The image that comes to mind --- semi-naked, forest-dwelling, illiterate, primitive people --- was \textbf{constructed by the European colonizers}.
\item In the 16th--18th centuries (the peak of colonization), categories such as \textbf{tribe, indigenous people, first nation, aborigine and native} were created to label the people of the non-Western world. The first label used was \textbf{native} (the original inhabitants of a place).
\end{itemize}
\end{KeyConcept}

\begin{ComparisonBox}
\begin{itemize}
\item In \textbf{Australia/America} there were indigenous inhabitants and external colonizers, so 'indigenous', 'first nation', 'native' all overlapped.
\item In \textbf{India}, everyone was native --- Brahmins, Kshatriyas, Vaishyas, Shudras and tribes alike --- so a fresh internal distinction had to be created: \textbf{caste vs tribe}.
\item The British divided Indian society into two categories: (1) those \textbf{absorbed in the mainstream feudal system} (the caste groups), and (2) those \textbf{relatively isolated, living in the forest} (the tribes).
\end{itemize}
\end{ComparisonBox}

\begin{KeyConcept}
\begin{itemize}
\item \textbf{The term 'tribe' was introduced by the British colonial government as an administrative strategy} --- to distinguish between those already incorporated within the hierarchical feudal system and those outside it.
\item This was a practice of \textbf{administrative convenience}: hierarchical (caste) society vs relatively egalitarian (tribal) society.
\end{itemize}
\end{KeyConcept}

\begin{itemize}
\item To create a distinct category, the colonizers used various \textbf{criteria}. The first was \textbf{religion}: those who practise \textbf{animism} (belief in soul / panpsychism --- a soul in every object) were considered tribes.
\item The criterion was later refined: those \textbf{outside the universal religions} (Hinduism, Islam, Sikhism, Christianity, Buddhism, Jainism) --- having their own indigenous/native religion --- would be considered tribes.
\end{itemize}

\subsection{The Classical Evolutionary Theory}

The definitions of tribe were intertwined with the political and social conditions of the time; two theories dominated the colonial era --- the classical evolutionary theory and functionalism.

\begin{KeyConcept}
\begin{itemize}
\item In the 16th century, the \textbf{Enlightenment} emerged in Europe, promoting \textbf{rational thought} over theological interpretation: society was no longer seen as divinely ordained, and every social relation came to be seen as a \textbf{human creation}.
\item With industrialisation, technology, inventions and navigation, Europeans travelled and met diverse cultures, and had to explain \textbf{human diversity}.
\item \textbf{Positivism} was established: society is external and objective, so data can be collected about it and \textbf{laws} of society (applying to every society) can be formulated.
\item This produced a \textbf{division of the social sciences}: \textbf{sociology} (to critically interpret one's own, industrial society) and \textbf{anthropology} (to study 'other cultures' --- the non-European, non-developed cultures).
\end{itemize}
\end{KeyConcept}

\begin{KeyConcept}
\begin{itemize}
\item Contact with non-Western peoples produced a \textbf{racial superiority complex}: Australian and African tribals exchanged gold for glass beads, so Europeans concluded these were 'child-like' people whose brains were not developed.
\item This formed the \textbf{evolutionary theory}: society develops in \textbf{stages}, and certain societies are \textbf{arrested} at a particular (early) stage of growth --- these were labelled \textbf{tribe} (least developed, uncivilised, illiterate, without good technology, chaotically arranged).
\item \textbf{Europeans speculated that these people represented the childhood or early stages of human development, and that Europeans had evolved into a superior stage they self-designated as 'civilisation'.} The word 'tribe' emerged in this context, applied to people seen as \textbf{far away from civilisation}.
\end{itemize}
\end{KeyConcept}

\begin{CriticalPoint}
\begin{itemize}
\item The animal-like equation of tribals allowed Europeans to \textbf{acquire their land and justify genocide} --- driving them away or killing them.
\item Tribal land was called \textbf{Terra Nullius} (terra = territory; nullius = void/zero) --- 'vacant land', ready to be captured, ignoring their habitation.
\item Most were killed in mass massacres; survivors were sent to \textbf{reserves} (like wildlife reserves/national parks) in America, Australia and New Zealand.
\item Example: the \textbf{Maori} of New Zealand --- the indigenous settlers subjugated to minority and vulnerability after colonisation (e.g. the wrestlers Dwayne 'The Rock' Johnson and Roman Reigns are Maori); recently a Maori candidate was in line to be the next Prime Minister of New Zealand.
\end{itemize}
\end{CriticalPoint}

\begin{GovtPolicy}
\begin{itemize}
\item \textbf{G. S. Ghurye} was inspired by evolutionism (and by \textbf{E. B. Tylor}, the founder of the school of evolutionism --- Ghurye even gave a racial origin of the caste system).
\item Ghurye defines tribes as communities which \textbf{live away from the civilised, speak the same tribal dialect, profess a primitive religion (animism), practise a primitive occupation (hunter-gatherers), are largely carnivorous, live naked or semi-naked, and have nomadic habits}.
\item The emphasis is on \textbf{primitiveness/backwardness} as a stage of evolution --- from which they will gradually develop towards civilisation.
\item This gave colonial powers the authority to rule: the \textbf{'White Man's Burden'} --- the colonisers claimed the obligation to civilise these 'uncivilised' people and share the fruits of civilisation. 'Civilisation' to them meant their own culture, the most evolved form.
\end{itemize}
\end{GovtPolicy}

\subsection{The Functionalist Approach}

The classical evolutionary theory was not based on fieldwork --- it relied on \textbf{second-hand data} (travellers, traders, missionaries, e.g. Marco Polo, Vasco da Gama), who sensationalised what they saw ('they eat raw meat, perform magic, roam naked, have incestuous sex'). This produced a distorted image, and the evolutionary theories were later rejected.

\begin{KeyConcept}
\begin{itemize}
\item \textbf{Functionalism} begins with \textbf{fieldwork} --- collecting \textbf{first-hand, empirical data} rather than speculation. The people among whom fieldwork was done were designated as \textbf{tribe}.
\item Functionalism put forth the paradigm of tribe as an \textbf{isolated, bounded and homogeneous entity}.
\item \textbf{Each tribe was a single homogeneous entity} --- a single pattern of institutions and ideas --- and those who shared this pattern participated in no other pattern. They were viewed as having a \textbf{definite identity}, endorsed by the practice of \textbf{endogamy} and a body of belief unique to themselves.
\end{itemize}
\end{KeyConcept}

\begin{itemize}
\item Rooted in \textbf{structural-functionalism}, this stressed \textbf{equilibrium}: the boundedness of social relations maintains equilibrium, so the tribe is presented as \textbf{unchanging, static, and ahistorical} --- living the same way for generations, with change coming only when they contact colonial powers.
\item \textbf{Boundedness} means one's whole life is defined by social relations --- especially \textbf{kinship and descent}. By birth into a family/descent group one gets fixed economic, political and social roles (no choices: e.g. the eldest son rules, the youngest cannot).
\end{itemize}

\begin{KeyConcept}
\begin{itemize}
\item \textbf{D. N. Majumdar} (an Indian structural-functionalist, inspired by Radcliffe-Brown) defines \textbf{tribe} as \textbf{a collection of families wearing a common name, occupying a common territory, observing certain taboos regarding marriage, profession/occupation, and having a well-assigned system of reciprocity}.
\item \textbf{Reciprocity} --- mutual obligation: giving something and expecting something in return. As \textbf{Marcel Mauss} (a relative of Durkheim, author of \textbf{The Gift}) put it, \textbf{'if friends make gifts, gifts make friends'} --- exchanging objects creates and maintains a social relationship, so strangers cannot simply transact as in a market.
\end{itemize}
\end{KeyConcept}

\begin{CriticalPoint}
\begin{itemize}
\item Functionalism was criticised on two grounds:
\item 1. \textbf{The neat classification was impossible} --- tribes cannot be neatly boxed as isolated/bounded. The \textbf{Ramayana} (from the Gupta age) shows \textbf{Shabri} (a tribal woman) and a \textbf{king of the forest} who helps Rama cross the river --- evidence of ancient contact between tribe and caste, thousands of years before colonisation.
\item 2. \textbf{Over-emphasis on equilibrium} --- the claim of a static, unchanging, ahistorical society was challenged (no society is without change/conflict).
\end{itemize}
\end{CriticalPoint}

\subsection{Elman Service's Socio-Political Typology}

\begin{KeyConcept}
\begin{itemize}
\item \textbf{Elman Service} used a \textbf{socio-political typology} to understand the tribe, dividing all the societies of the world, on the basis of organisational complexity, into \textbf{four categories}: \textbf{band, tribe, chiefdom, state}.
\item The \textbf{tribe is simply a stage in organisational complexity} --- more complex than a band, less complex than a chiefdom and state --- with \textbf{no developmental stigma or progressive ideology attached}.
\item \textbf{Tribe refers to all those societies based on relations of kinship, especially descent groups}: a person becomes a member by being \textbf{born into a descent group}, and carries that identity for purposes of \textbf{marriage, ownership of resources, and political relations}.
\item This is the more comprehensive definition accepted by scholars --- it places tribe on a scale of evolving organisational complexity without attaching any developmental or progressive significance to the stages.
\end{itemize}
\end{KeyConcept}

\subsection{Post-Colonial Paradigm Shift}

Colonial ethnography was sponsored by the colonial rulers and could not criticise colonial rule --- it painted a happy-go-lucky picture of tribes as isolated societies in equilibrium, ignoring the violence, genocide and epidemics that colonial rule brought.

\begin{KeyConcept}
\begin{itemize}
\item In the \textbf{post-colonial period}, when natives got the chance to write about themselves, scholars began to project tribes as a \textbf{distinct ethnic group} --- not racial, not backward, not static.
\item An \textbf{ethnic group} has: (1) a \textbf{sense of unity / identity consciousness} (I belong to that group), and (2) a \textbf{shared experience} of culture, history, or another unifying factor.
\item \textbf{Post-colonial ethnography rejected the idea of equilibrium / static tribal society, and the term 'tribe' was replaced by 'ethnic groups'} --- applied mostly to communities that may be \textbf{less dominant in the social hierarchy} but have a \textbf{distinct culture}.
\item Tribes are now seen as \textbf{interacting} with the mainstream and \textbf{gradually transforming} --- e.g. Scheduled Tribe members entering the IAS/IPS and white-collar jobs, and tribal tourism (museums, restaurants, souvenirs, the 'tribal circuit') showcasing their culture. They can no longer be labelled backward, primitive, or static.
\end{itemize}
\end{KeyConcept}

The debate on whether \textbf{tribes should be seen as the 'indigenous people of India'} (an EPW article by the sociologist \textbf{Virginius Xaxa}) is taken up in the next class.

\begin{QuickRevision}
\begin{itemize}
\item Four themes: definitional problem, geographical spread, colonial policies, integration \& autonomy --- all rooted in colonial history.
\item 'Tribe' is an external, colonially imposed category --- tribals do not call themselves tribes (Mundas = 'Horoko', original creation of creator; S. C. Roy).
\item British administrative strategy: caste = absorbed in feudal hierarchy; tribe = isolated forest-dwellers; criteria --- animism, outside universal religions.
\item Classical evolutionary theory: Enlightenment $\rightarrow$ rationality $\rightarrow$ positivism $\rightarrow$ racial superiority complex $\rightarrow$ tribe = 'arrested childhood stage'; terra nullius; genocide/reserves; Ghurye's definition; White Man's Burden.
\item Functionalism: fieldwork \& first-hand data; tribe = isolated, bounded, homogeneous, static (equilibrium); D. N. Majumdar's definition; reciprocity (Mauss, 'The Gift'); criticised (Ramayana's Shabri; no pure equilibrium).
\item Elman Service: band $\rightarrow$ tribe $\rightarrow$ chiefdom $\rightarrow$ state; tribe = kinship/descent-based, no stigma.
\item Post-colonial shift: 'tribe' $\rightarrow$ 'ethnic groups' (identity consciousness + shared culture), interacting and transforming.
\end{itemize}
\end{QuickRevision}

\section{Indigenous Status and the Tribe--Caste Relationship}

\subsection{Who Are 'Indigenous' People?}

In lay terms, \textbf{indigenous} means the original inhabitant of a place --- but today the relevant idea is that they are the \textbf{descendants} of original inhabitants, not the originals themselves.

\begin{GovtPolicy}
\begin{itemize}
\item The term gained currency through the \textbf{International Labour Organization (ILO)}.
\item \textbf{ILO Convention 107 (1957)} defined who is indigenous.
\item \textbf{ILO Convention 169 (1989)} gave the idea of how to \emph{identify} indigenous people.
\item ILO's position: it is more important to \textbf{identify} indigenous people than to \textbf{define} them.
\end{itemize}
\end{GovtPolicy}

\begin{KeyConcept}
\begin{itemize}
\item ILO's identification criteria:
\item 1. \textbf{Self-identification / self-determination} --- they proclaim themselves the original inhabitants (and we accept it).
\item 2. \textbf{Historical continuity} --- a historical connection with the ancestors whose descendants they claim to be.
\item 3. \textbf{Distinct political, social and economic system} --- distinct from the mainstream, having maintained their distinct identity and not fully assimilated.
\item 4. \textbf{Distinct language, culture and belief} --- distinct way of life, thought process, religion, family and other institutions.
\item 5. \textbf{Non-dominant part of society} --- i.e. marginalised.
\item In short: original inhabitants before colonisation, with historical continuity, distinct systems and culture, and marginalised status.
\end{itemize}
\end{KeyConcept}

\subsection{The Indigenous-Status Debate: Virginius Xaxa}

The debate over whether Indian tribes are indigenous is summarised from an EPW article by the sociologist \textbf{Virginius Xaxa} ('Tribes as Indigenous People of India'). Until 1989 tribes were accepted as indigenous; after 1989 the Indian government had difficulty accepting the label.

\begin{KeyConcept}
\begin{itemize}
\item The definition of tribe has always been difficult, but one thing associated with tribes is accepted: they are the \textbf{dispossessed and deprived people of a region} --- vulnerable, marginalised, needing support.
\item Indian tribes have \textbf{never claimed to be the original inhabitants}; what they have asked for is \textbf{access to the resources of their region} and a \textbf{say in how those resources are used}.
\item Tribes occupy regions rich in forest and natural resources (minerals, coal), but these resources have been extracted \textbf{without sharing the fruits of development} --- their land is snatched by corporates or government companies and they are displaced.
\item The tribal identity gives self-esteem and has been used to mobilise people for their rights and due share of development --- in the process they have begun to assert themselves as the indigenous population/natives of India.
\end{itemize}
\end{KeyConcept}

\subsection{ILO Conventions 107 and 169}

\begin{GovtPolicy}
\begin{itemize}
\item \textbf{Convention 107 (1957)} --- about the integration of indigenous people: to 'protect and promote the progressive integration' of indigenous people into mainstream society. India \textbf{accepted} this convention.
\item \textbf{Integration $\neq$ assimilation}: assimilation means losing one's identity and becoming part of the other; integration means keeping one's identity while becoming part of the whole --- like a \textbf{rainbow} (each colour keeps its identity yet is part of the continuum) or a \textbf{salad bowl}.
\item Example: the \textbf{Sikh} community has not become Hindu to integrate into India --- it keeps its own identity, religion, and separate laws, yet is part of the nation.
\end{itemize}
\end{GovtPolicy}

\begin{GovtPolicy}
\begin{itemize}
\item \textbf{1985} --- on the basis of several activists, ILO concluded the convention needed revision: integration vs assimilation should be \textbf{decided by the tribes themselves} (a 'next level of democracy' --- they should plan and execute their own development schemes).
\item \textbf{Convention 169 (1989)} --- shifts from \textbf{integration to empowerment}: it speaks of \textbf{need rights} and \textbf{power rights}, especially \textbf{rights over land} (what is to be done with their land should be their decision, not the government's or a mining company's).
\item Convention 169 gave a specific definition of indigenous highlighting \textbf{three aspects}: (1) \textbf{lived before colonisation}; (2) \textbf{marginalised after colonisation}; (3) \textbf{govern their life based on their own socio-economic and cultural institutions} (rather than the broad constitutional/state structure).
\item This created a problem for the Indian government, which therefore \textbf{did not accept Convention 169}.
\end{itemize}
\end{GovtPolicy}

\subsection{Arguments Against Indigenous Status}

\begin{CriticalPoint}
\begin{itemize}
\item The government's core objection: \textbf{no population in India is less indigenous than the tribes} --- Sikhs, Hindus, Muslims, Buddhists and Jains were all present before colonisation, so none is 'more indigenous' than another.
\item The detailed arguments against granting tribes indigenous status:
\end{itemize}
\end{CriticalPoint}

\begin{itemize}
\item \textbf{Waves, not one wave, of colonisation.} Unlike the US, New Zealand and Australia (where a single wave of colonisation in the 16th--17th century pushed the original population to marginal areas), India has seen waves of movement over centuries and millennia --- so there is no fixed point to decide who is native. \textbf{How far back should one go to determine who is native and who is not?}
\item \textbf{Close proximity and acculturation.} Tribes have lived in close proximity with non-tribes for centuries, leading to much acculturation and even assimilation (e.g. the \textbf{Ramayana}'s Bhil woman \textbf{Shabri}, who offers berries/fruits to Rama --- evidence of smooth cultural contact). Indian experience is different from the US pattern of pushing natives to reserves.
\item \textbf{The Aryan-migration problem.} If the Aryan migration (which pushed the original inhabitants southward) is taken as the reference, then non-tribes become the 'migrants'. But many tribes settled \emph{after} the Aryans: the \textbf{Mongoloid tribes of the North-East} migrated from Tibet, Myanmar and Mongolia (via China) and settled only around the 1st millennium AD --- the \textbf{Meiteis} around the 16th century, the \textbf{Kukis} even later. So \textbf{not all tribes can claim settlement prior to the Aryans}.
\item \textbf{The reference-point problem.} Should indigenous status be given at the national or the local (territory-occupied) level? If \textbf{original settlement} is the criterion, several \textbf{non-tribal} communities would be included as indigenous --- the \textbf{Dravidian} speakers (Tamil, Kannada, Telugu, Malayalam) settled before the Aryans --- while several tribes (Nagas, Kukis, Mizos) would be excluded. These non-tribes are also often the \emph{dominant} communities of their states, so they cannot be given indigenous status.
\item \textbf{Fused, mixed communities.} Indian society is made up of caste groups formed out of the \textbf{fusion} of various groups --- Bengalis, Odiyas and others have both indigenous and non-indigenous components, so only \textbf{one segment} of a community would be 'indigenous' while the rest would not.
\end{itemize}

\subsection{Arguments For Indigenous Status}

\begin{KeyConcept}
\begin{itemize}
\item Activists, tribal scholars and sociologists (including \textbf{Virginius Xaxa}) argue for declaring tribes indigenous because it leads to their \textbf{self-empowerment} --- improved welfare and participatory democracy (they plan and execute their own welfare plans).
\item The claim should not be read literally (original settlement) but for the \textbf{larger good} it brings.
\end{itemize}
\end{KeyConcept}

\begin{itemize}
\item \textbf{Time of settlement, not original settlement.} It is defended not on the basis of \emph{original} settlement but \emph{time} of settlement --- their fathers, grandfathers and forefathers have lived in that area for a very long time, so they have original-settler status of that place.
\item \textbf{Worst victims of colonisation.} Forests were exploited for raw material, tribals were treated like animals and their land treated as \textbf{terra nullius} --- so they faced the maximum brunt of colonisation and share all the \textbf{attributes of colonised people} (an ethnic identity as a non-dominant, vulnerable section).
\item \textbf{Do not mix the two levels.} Do not confuse the national level with the local level --- simply give them original-settler/priority rights over the particular area they have long inhabited.
\item \textbf{Understand why the claim is arising.} Dominant communities make the same 'this is our land' claim through \textbf{sons-of-the-soil movements} --- Marathi speakers in Maharashtra (against 'Biharis', during Balasaheb Thackeray and Raj Thackeray's agitations), Tamils in Tamil Nadu, Kannadigas, Punjabis, Gujaratis --- and governments respond with reservations (e.g. Maharashtra/Haryana reserving private jobs for locals). But this privilege is denied to tribals: coal and minerals are extracted on their land (by Vedanta, Tata, etc.) without them getting the benefits or the right to protest. The \textbf{minimum} that should be done is to \textbf{recognise their prior rights over their land and resources}.
\end{itemize}

\begin{QuickRevision}
\begin{itemize}
\item Indigenous = descendants of original inhabitants; ILO 107 (1957) defined it, ILO 169 (1989) identified it (self-identification, historical continuity, distinct systems/culture, non-dominant).
\item Xaxa's debate: tribes are dispossessed/deprived; they never claimed to be 'original' --- they demand rights over regional resources.
\item 107 = progressive integration (rainbow/salad bowl, not assimilation); India accepted it. 169 = empowerment + land rights; India rejected it.
\item Against indigenous status: waves of migration (no fixed point); proximity/acculturation (Shabri); Aryan migration (NE tribes settled later); reference point (Dravidians would count, Nagas/Kukis/Mizos would not); fused communities (one segment only).
\item For indigenous status: time (not original) of settlement; worst victims of colonisation (terra nullius); don't mix levels; sons-of-the-soil hypocrisy --- minimum is recognising land/resource rights.
\end{itemize}
\end{QuickRevision}

\subsection{Tribe and Caste: Two Foreign Categories}

\begin{KeyConcept}
\begin{itemize}
\item Both 'tribe' and 'caste' are \textbf{foreign words, not Indian} --- neither is how Indian society organises itself.
\item \textbf{Tribe} is a colonial construct (introduced by the British).
\item \textbf{Caste} comes from the Portuguese word \textbf{'casta'} --- the Portuguese saw Indian society divided into distinct, endogamous groups and used 'caste' to identify them.
\item The Indian terms: tribe = \textbf{janjati}; caste = \textbf{jati} (the 'ground reality'). Janjati is simply the \textbf{jati that lives in the forest}.
\item Colonial ethnographers wanted to create a \textbf{watertight boundary} between caste and tribe --- projecting tribe as a distinct, isolated group unaffected by caste.
\end{itemize}
\end{KeyConcept}

Post-colonisation, scholars showed that tribe and caste were never separate --- they interacted, and their identities are \textbf{fluid, not fixed}. The evolutionists defined tribe as backward; the functionalists as isolated/bounded with a unique fixed identity. This categorisation became official through the \textbf{Census} (1901: those following animism; 1931: those following a religion other than the mainstream).

\subsection{From Watertight Boundary to Continuum: Mandelbaum and Wiser}

\begin{KeyConcept}
\begin{itemize}
\item \textbf{Mandelbaum} (David Mandelbaum) --- there are many groups at the \textbf{intersection of tribe and caste} whose identities are difficult to pinpoint (they appear like both a tribe and a caste).
\item His distinguishing criterion: \textbf{in a tribe, the principal links for the whole of society are based on kinship} --- but for a \textbf{jati}, it is not just kinship but the \textbf{presence of other jatis of unequal rank} that enables it to continue as a particular jati group.
\item Therefore \textbf{hierarchy is the central feature of caste} (absent in the egalitarian tribe): jatis insist on hierarchical ordering, while tribesmen organise society by kinship without insisting on hierarchy.
\end{itemize}
\end{KeyConcept}

\begin{ExampleBox}
\textbf{William Wiser} --- in the village, several jatis performing different tasks live together in \textbf{organic solidarity} (interdependence due to heterogeneity: a barber cuts hair, a carpenter offers services). This heterogeneity is \textbf{not} the case with tribes, where homogeneity prevails.
\end{ExampleBox}

\subsection{The Tribe--Caste Continuum}

\begin{KeyConcept}
\begin{itemize}
\item A \textbf{continuum} means there is no discrete categorisation --- change from one end to the other is so gradual that you cannot draw a clean line (like a rainbow: you cannot tell exactly where violet ends and indigo begins).
\item The \textbf{tribe--caste continuum} holds that \textbf{tribe and caste are polar ideal categories} --- ideal types (in Weber's sense) that do not exist in reality. They are reference points; most communities lie \textbf{somewhere in between} --- closer to tribe but with some caste influence, or closer to caste but with some tribal influence.
\item This idea was given by \textbf{F. G. Bailey} (who studied the village \textbf{Bishipara in Odisha}), \textbf{Surjit Sinha} (who studied the \textbf{Lodha} community in West Bengal), and \textbf{D. N. Majumdar}.
\item The continuum reflects continuous interaction: tribes and castes have influenced each other and adopted elements from each other --- no tribe is completely isolated from caste influence.
\end{itemize}
\end{KeyConcept}

\begin{ExampleBox}
\begin{itemize}
\item \textbf{D. N. Majumdar} --- the movement of tribal people adopting caste elements is \textbf{one form of Sanskritization} (using M. N. Srinivas's idea): tribes absorb caste-based ideology/elements to be recognised in Hindu society and proclaim a caste identity.
\item Example from the \textbf{Garo Hills}: the \textbf{Songsarek Garos} show little or no sign of Sanskritization, while the \textbf{Dalus} have completely adopted a Hindu identity (proclaiming themselves Vaishnavas). These are the two extremes of the continuum; in between lie \textbf{Rabha, Koch and Hajong} --- absorbing caste elements but not fully identifying as caste.
\item Christian missionaries established schools and converted tribals (especially in Jharkhand, Chhattisgarh, Odisha and the North-East); later organisations like the \textbf{RSS} began similar activity (e.g. coaching institutes in Arunachal) inculcating Hindu values --- both processes erode tribal identity.
\end{itemize}
\end{ExampleBox}

\subsection{Hindu Mode of Tribal Absorption (N. K. Bose)}

\begin{KeyConcept}
\begin{itemize}
\item \textbf{Nirmal Kumar Bose (N. K. Bose)} gave the idea of the \textbf{'Hindu mode of tribal absorption'} --- the process by which Hindus gradually absorb a tribal community as a distinct caste group.
\item Reasons why a tribe adopts the caste system:
\item 1. \textbf{Monopoly over an occupation} --- economic security (their basic economy runs; they are no longer dependent).
\item 2. The \textbf{level of technology and lifestyle} the tribe saw among Hindu farmers/peasants.
\item A tribe absorbs caste identity \textbf{irrespective of the position it gets in the hierarchy}.
\end{itemize}
\end{KeyConcept}

\begin{ExampleBox}
\begin{itemize}
\item Example: the \textbf{Juangs of Odisha} (studied by Bose) --- the dominant castes depended on the Juangs for local liquor, became addicted, and fell into debt; through these economic ties with mainstream castes, the Juangs gradually assumed a caste-like identity.
\item \textbf{Poverty of Brahmins} was another reason (also cited by \textbf{Herbert Risley}): poor priests acted as guides for aspiring tribes --- for a fee they created stories proving a tribe's claim to a dominant (e.g. Rajput) caste identity.
\item Example: the \textbf{Kanjar} community --- said to be Rajputs and descendants of \textbf{Maharana Pratap} who left their kingdom after the Mughal/Delhi Sultanate invasion and took to the forest, thus acquiring a tribal identity. Such stories legitimised the tribe's upward claim.
\end{itemize}
\end{ExampleBox}

\begin{CriticalPoint}
\textbf{N. K. Bose is criticised} for talking about movement in \textbf{one direction only} (tribe $\rightarrow$ caste). In fact movement also happens the other way --- \textbf{caste moving towards tribe} --- which is called \textbf{tribalization}.
\end{CriticalPoint}

\subsection{Tribalization: The Reverse Movement (C. W. Brown)}

\begin{ExampleBox}
\begin{itemize}
\item \textbf{C. W. Brown} studied the \textbf{Bhotias} (now in Uttarakhand, at the Indo-Tibetan border): they were actually \textbf{Rajputs} who preferred a \textbf{non-Hindu lifestyle} (e.g. eating beef) in order to form long-term trading relations with Tibetans (which involved living in each other's houses).
\item They \textbf{sacrificed their upper-caste status for economic benefit} --- becoming nomadic traders with a tribal-like identity.
\item \textbf{Post-independence}, cross-border trade was discontinued, and the Bhotias --- whose economic and political relations were now within mainland India --- \textbf{attempted to move back towards Hindu upper-caste identity}.
\item This shows movement from \textbf{caste to tribe} (tribalization), not just tribe to caste.
\end{itemize}
\end{ExampleBox}

\begin{KeyConcept}
\begin{itemize}
\item \textbf{Conclusion}: colonial scholars and the functionalist definition tried to project a \textbf{fixed} identity, but the identity of tribe and caste is in reality \textbf{fluid} --- it changes with context and circumstances (a group may move towards tribal identity or towards caste identity depending on the situation).
\item The identity of being a tribe or a caste is \textbf{a matter of historical and political process}, and \textbf{may change over time depending on the context}.
\end{itemize}
\end{KeyConcept}

\begin{CriticalPoint}
\textbf{Note on terminology}: Sanskritization and 'casteization' are \textbf{not interchangeable}. Majumdar uses Srinivas's idea of Sanskritization to explain tribe$\rightarrow$caste mobility; it is \textbf{one type of Sanskritization}, not the same as casteization.
\end{CriticalPoint}

\begin{QuickRevision}
\begin{itemize}
\item 'Tribe' and 'caste' are foreign constructs --- tribe (colonial), caste (Portuguese 'casta'); Indian terms are \textbf{janjati} (forest jati) and \textbf{jati}.
\item Colonial ethnographers drew a watertight caste--tribe boundary; post-colonial scholars show it was fluid and interactive.
\item Mandelbaum: tribe = kinship-based, egalitarian; jati = hierarchy of unequal-ranked jatis. Wiser: village = heterogeneous organic solidarity.
\item Tribe--caste continuum (Bailey--Bishipara, Sinha--Lodha, Majumdar): polar ideal types, most groups in between; Majumdar = tribe$\rightarrow$caste as a form of Sanskritization (Garo Hills: Songsarek $\leftrightarrow$ Dalus, with Rabha/Koch/Hajong in between).
\item N. K. Bose --- Hindu mode of tribal absorption (occupation monopoly, technology/lifestyle, Juangs, poverty of Brahmins/Risley, Kanjar--Maharana Pratap story); criticised for one-direction movement.
\item Tribalization (reverse): C. W. Brown's Bhotias (Rajputs who ate beef for Tibetan trade).
\item Tribe/caste identity = historical-political process, changes with context.
\end{itemize}
\end{QuickRevision}

\section{Scheduled Tribes, PVTGs and Geographical Spread}

\subsection{From Tribe to Scheduled Tribe}

The journey is from a socio-cultural entity to a politico-administrative category. The term \textbf{tribe} originated with British rule (from 1885), which identified a feudal society as \textbf{caste} and a more backward one as \textbf{tribe}; later functionalism pictured the tribe as a close, bonded, homogeneous society tied to a territory --- so the colonial idea of the tribe was \textbf{backward} and \textbf{separated}.

\begin{KeyConcept}
\begin{itemize}
\item \textbf{Tribes were identified largely in terms of what they were not} --- not practising Vedic Hinduism, not Muslim, no stratification, no money economy. These were \textbf{negative traits} (negation/absence of a characteristic), not positive ones.
\item The \textbf{Census of India} (started by the British) consolidated this identity:
\item \textbf{1901} --- tribes = those who follow \textbf{animism}.
\item \textbf{1921} --- the additional idea of \textbf{territory} (localised people with a specific identity).
\item \textbf{1931} --- the term \textbf{'primitive tribe'} was introduced.
\item \textbf{Government of India Act 1935} --- 'primitive tribe' was replaced by \textbf{'backward tribe'}.
\item \textbf{1941} --- tribes identified by \textbf{tribal origin} (subjective, giving the government flexibility to classify anyone as tribe or caste).
\end{itemize}
\end{KeyConcept}

Post-independence scholars challenged this: the colonial ethnography had rendered a \textbf{fixed image} to an identity that was actually \textbf{fluid and contextual} (with instances of both tribalization and casteization).

\begin{KeyConcept}
\begin{itemize}
\item \textbf{Tribe} = a \textbf{cultural category} (culturally fluid, no single criterion separates it from caste).
\item \textbf{Scheduled Tribe (ST)} = a \textbf{politico-administrative category} --- objective (you are or you are not), with no cultural elements. It does \textbf{not capture the enormous socio-cultural complexity} of the various tribes.
\item The categorisation of the population was done \textbf{to make special welfare provision available} to them.
\end{itemize}
\end{KeyConcept}

\begin{GovtPolicy}
\begin{itemize}
\item Two key articles: \textbf{Article 366} (defines Scheduled Tribes as those deemed to be so under Article 342) and \textbf{Article 342} (provides for the listing of Scheduled Tribes).
\item The same groups listed under the Government of India Act 1935 were reclassified as \textbf{Scheduled Tribes in 1950}, with constitutional/legal provisions (reservations and other rights) attached.
\end{itemize}
\end{GovtPolicy}

\subsection{Ghurye vs Elwin: Assimilation vs Isolation}

\begin{ComparisonBox}
\begin{itemize}
\item \textbf{G. S. Ghurye --- assimilation.} An Indologist (famous for the racial origin of caste), he held that tribes are \textbf{'backward Hindus'} --- not different from caste --- and the best way to manage them is to \textbf{assimilate} them into the mainstream Hindu value system, breaking down all barriers.
\item \textbf{Verrier Elwin --- isolation.} Based on his fieldwork, he held that whenever tribes come into contact with the mainstream they are \textbf{exploited} and their resources snatched --- so they should be \textbf{protected and isolated}, like a \textbf{national reserve / wildlife sanctuary} (minimal contact with outsiders).
\item Elwin was criticised as a \textbf{'colonial hangover'} (the US pattern of pushing indigenous people to reserves, making them tenants in their own land).
\item The Indian government therefore chose a middle path: \textbf{Nehru's integration approach} (the \textbf{tribal panchsheel} framework) --- integrate, neither isolate nor assimilate.
\end{itemize}
\end{ComparisonBox}

\subsection{Identifying the Scheduled Tribe: Commissions and the Lokur Criteria}

The initial criterion was borrowed from the Government of India Act 1935. Later, the \textbf{First Backward Classes Commission (Kaka Kalelkar)} made some suggestions (not accepted for classifying tribes), and in \textbf{1965} the \textbf{Lokur Committee} gave specific criteria.

\begin{GovtPolicy}
\begin{itemize}
\item \textbf{Lokur Committee (1965)} --- five criteria to identify tribes:
\item 1. \textbf{Primitive traits}.
\item 2. \textbf{Geographical isolation}.
\item 3. \textbf{Shyness of contact with the community at large}.
\item 4. \textbf{Distinct culture}.
\item 5. \textbf{Backwardness}.
\end{itemize}
\end{GovtPolicy}

\begin{CriticalPoint}
\begin{itemize}
\item These criteria have \textbf{not been revised since 1965} and are now largely \textbf{obsolete}:
\item \textbf{Primitive traits} --- no objective definition of 'primitive'; many tribes have modernised and integrated with the mainstream.
\item \textbf{Geographical isolation} --- in a globalised world (trains, ships, tribal circuits, tourism in Andaman/Lakshadweep) no part of India is isolated any more.
\item \textbf{Shyness of contact} --- largely obsolete; ST students study and interact with everyone. The \textbf{Sentinelese} (who killed an American missionary trying to enter their island) are a rare exception.
\item \textbf{Distinct culture} --- the idea rested on a flawed assumption of a single Brahminical-Hindu 'national culture'; India is diverse from Kashmir to Kanyakumari, yet with common threads that are increasing with globalisation, internet and media.
\item \textbf{Backwardness} --- again subjective (economic? educational?), with no clear criterion.
\item The process of inclusion/exclusion is therefore \textbf{vague and based on flawed assumptions}.
\end{itemize}
\end{CriticalPoint}

\subsection{Area Restriction and Its Removal (1976)}

\begin{KeyConcept}
\begin{itemize}
\item Under the British-era criteria, tribes were assumed to be \textbf{localised groups}, so tribal status applied only within a \textbf{specific territory}: a person could be a Scheduled Tribe in one district but not in the neighbouring district.
\item This \textbf{area restriction} discouraged mobility --- a tribe moving to another state for work would lose the welfare benefits.
\item The \textbf{Lokur Committee} argued the territorial restriction was a \textbf{barrier to spatial and social mobility} and \textbf{contrary to the goal of tribal integration}, and said it should be removed.
\item \textbf{1976 --- Removal of Area Restrictions Act}: restriction lifted, but now at the \textbf{state level} (a tribe is recognised in a particular state, not nationally) --- e.g. the \textbf{Naga} are STs in Nagaland, not in UP; the \textbf{Banjara} are STs in several states but included in the SC list in Delhi (which has no ST list).
\end{itemize}
\end{KeyConcept}

\subsection{Anomalies in Scheduling and the Political Gimmick}

\begin{CriticalPoint}
\begin{itemize}
\item With no objective criteria for inclusion/exclusion, scheduling becomes a \textbf{political gimmick/tool} to mobilise people and gather votes --- e.g. \textbf{Gurjars} and other communities demanding ST status.
\item \textbf{Anomalies}: some communities that have improved and integrated do not want to give up their ST identity (ST is politico-administrative, not cultural); and casteist people wrongly treat STs as a lower caste, though ST communities lie \textbf{outside the caste system}.
\item These anomalies need a commission to be corrected.
\end{itemize}
\end{CriticalPoint}

\begin{ExampleBox}
\begin{itemize}
\item Related reference (from \textbf{Christophe Jaffrelot}, the caste specialist, in Indian Express): why dominant castes like \textbf{Patidar, Patel and Maratha} now demand \textbf{OBC} status.
\item The \textbf{Green Revolution} produced a rising middle caste whose dominance rested on \textbf{land} and \textbf{economic superiority} (agriculture). But agriculture has lost its remunerative value, and land is no longer a source of dominance --- wealth now comes from \textbf{private jobs, IT and startups} (captured by educated Brahmins/Kshatriyas) and government jobs (where SCs benefit from reservation).
\item Left only with \textbf{numerical preponderance}, these middle castes demand OBC status (and even sub-quotas within OBC) to regain status and compete for government jobs.
\end{itemize}
\end{ExampleBox}

\begin{QuickRevision}
\begin{itemize}
\item Tribe (cultural) $\rightarrow$ Scheduled Tribe (politico-administrative) via colonial census: 1901 animism, 1921 territory, 1931 'primitive', 1935 'backward', 1941 origin.
\item Tribes identified by negative traits ('what they were not').
\item Ghurye = assimilation (backward Hindus); Elwin = isolation (national-reserve model, 'colonial hangover'); Nehru = integration (tribal panchsheel).
\item Articles 366 \& 342; Kaka Kalelkar commission; Lokur Committee 1965 (five criteria).
\item Area restriction removed 1976 $\rightarrow$ state-level recognition.
\item Criteria obsolete; scheduling a political gimmick; anomalies persist.
\end{itemize}
\end{QuickRevision}

\subsection{PVTGs: Origin and Criteria}

\begin{GovtPolicy}
\begin{itemize}
\item The \textbf{Dhebar Commission} created \textbf{PTG (Primitive Tribal Group)} as a separate category in \textbf{1973}, out of all the Scheduled Tribes.
\item In \textbf{2006}, PTG was renamed \textbf{PVTG (Particularly Vulnerable Tribal Group)} --- dropping the derogatory term 'primitive'.
\item PVTGs are tribes that are \textbf{more marginalised and vulnerable even among the STs} --- on the verge of extinction if not specially protected.
\end{itemize}
\end{GovtPolicy}

\begin{DataFact}
\begin{itemize}
\item Currently \textbf{75 groups} are listed as PVTGs; the maximum (\textbf{13}) are in \textbf{Odisha}.
\item They are found mainly in \textbf{seven states}: Maharashtra, Madhya Pradesh, Chhattisgarh, Jharkhand, Tamil Nadu, Andhra Pradesh and Odisha.
\item Total PVTG population $\approx$ \textbf{27.5 lakh}.
\item Famous PVTGs: \textbf{Jarawas} and \textbf{Sentinelese} (Andaman), and the \textbf{Bonda} (Odisha).
\end{itemize}
\end{DataFact}

\begin{KeyConcept}
\begin{itemize}
\item The \textbf{criteria} for a PVTG:
\item 1. \textbf{Forest-dependent livelihood} (living in the forest).
\item 2. \textbf{Pre-agricultural level of existence} --- a way of life similar to the Stone Age (Paleolithic/Mesolithic), before the Neolithic agricultural stage.
\item 3. \textbf{Stagnant or declining population}.
\item 4. \textbf{Low literacy rate}.
\item 5. \textbf{Subsistence-based economy} --- no concept of surplus (they produce only for consumption, unlike a modern surplus/profit economy).
\end{itemize}
\end{KeyConcept}

\begin{ExampleBox}
\begin{itemize}
\item \textbf{Mankidia} --- a community living in the \textbf{Simlipal Tiger Reserve} (Odisha) who used to kill monkeys for food; they were excluded/removed when the reserve was created.
\item \textbf{Jarawas} --- the \textbf{Andaman Trunk Road} passes through their habitation, and 'human safari' tourism treats them as a showpiece (people photograph them); police officers were once suspended for giving Jarawa girls biscuits to make them dance naked.
\item \textbf{Bonda (Odisha)} --- Bonda women wear only a two-inch cloth called \textbf{ringa} to cover the lower abdomen (with beads over the chest).
\item The \textbf{folklore} (narrated by \textbf{N. K. Bose} and \textbf{Verrier Elwin}): when Sita was bathing in a lake in Odisha, Bonda women stole her saree; on discovering this she cursed them to remain naked; when they begged forgiveness she tore a piece of her saree and gave it to them --- hence the divine instruction to wear the ringa. Two educated Bonda girls who rejected the custom were labelled witches and physically punished.
\item This shows how stories are created to absorb communities into the Hindu fold --- e.g. \textbf{Ganesha} (an elephant totem of a tribal community turned into a god) and \textbf{Kali} (a tribal goddess depicted as an incarnation of Durga).
\end{itemize}
\end{ExampleBox}

\subsection{Problems Faced by PVTGs}

\begin{CriticalPoint}
\begin{itemize}
\item 1. \textbf{Loss of traditional livelihood and customary rights on resources} --- industrial projects, conservation and tourism have removed them from forest land (the Mankidia from Simlipal; they are the \textbf{'refugees of development'}). Removed from the forest they lose their only skill (exploiting forest resources --- fuel, food, timber, medicine) and face a deep trauma (like being shifted from a palatial village home to a 2BHK flat).
\item 2. \textbf{Perceived primitivism/backwardness in discourse} --- the term 'primitive tribal group' is still used in official communication (even in 2021), devaluing their culture and tradition.
\item 3. \textbf{Irrational criteria of inclusion} --- inclusion/exclusion is \textbf{ad hoc} with no logic (some Odisha STs satisfy all criteria but are not listed), so many communities miss special assistance.
\item 4. \textbf{Treated as 'endangered'} --- socio-economic vulnerability led them to be treated as on the verge of extinction, leading to \textbf{disastrous intervention}: disallowing PVTG members from availing \textbf{sterilisation schemes} (their reproductive rights are not recognised). This \textbf{sidesteps the real factors for mortality} (poor nutrition, ineffective vaccination, lack of healthcare). A recent Chhattisgarh High Court judgment and an MP case relaxed this, but they still need SDM permission for reproductive rights.
\item 5. \textbf{Bonded labour} --- when their land is alienated and they are forced out, they end up as manual/bonded labourers (e.g. \textbf{Sahariyas} of Rajasthan, \textbf{Juangs} of Odisha) through systems like the \textbf{Haliya} system of debt-bondage, where they cannot even pay the interest from their wages.
\item 6. \textbf{Low literacy} --- in single digits, due to poor education policy, infrastructure shortage and teacher absenteeism in PVTG areas.
\item 7. \textbf{Loss of land ownership} --- e.g. 245 \textbf{Baiga} families removed for the \textbf{Achanakmar Tiger Reserve}.
\item 8. \textbf{Lack of baseline surveys} --- an Anthropological Survey of India report notes there is no vital data (infant/maternal mortality, education, per-capita income) for most PVTGs, so no targets can be set for schemes --- reflecting government ignorance.
\end{itemize}
\end{CriticalPoint}

\subsection{Scheme and Recommendations (NAC 2013)}

\begin{GovtPolicy}
\begin{itemize}
\item \textbf{Scheme for Development of PVTGs} --- the state government plans the scheme (education, health, economic livelihood) with full flexibility and submits it to the Tribal Ministry; the \textbf{central government provides 100\% funding}.
\item Problem: different PVTGs have different problems, so a \textbf{'one size fits all'} approach cannot work --- the scheme remains largely \textbf{top-down}.
\item The \textbf{National Advisory Committee (NAC, 2013)} --- headed by Sonia Gandhi --- recommended:
\item 1. \textbf{Status of Scheduled Tribe to all PVTGs} (some PVTGs, e.g. the \textbf{Paudi Bhuinya}, are not even recognised as STs).
\item 2. A \textbf{rights-based approach} to empowerment --- ensuring \textbf{habitat rights}, land rights, customary habitation, and educational/economic rights (restoring their habitation on forest land).
\item 3. A \textbf{single-window approach} --- one portal for all benefits (education, health, infrastructure, economy), instead of multiple windows.
\item 4. \textbf{Regular census and baseline surveys} so schemes can be customised.
\end{itemize}
\end{GovtPolicy}

\begin{QuickRevision}
\begin{itemize}
\item PTG (1973, Dhebar Commission) $\rightarrow$ PVTG (2006); 75 groups, max in Odisha (13); 7 states; \textasciitilde{}27.5 lakh.
\item PVTG criteria: forest-dependent livelihood; pre-agricultural (Stone Age) existence; stagnant/declining population; low literacy; subsistence economy (no surplus).
\item Examples: Mankidia (Simlipal), Jarawa/Sentinelese (Andaman), Bonda (ringa + Sita folklore).
\item Problems: loss of livelihood/rights ('refugees of development'); perceived primitivism; ad-hoc inclusion; reproductive-rights denial; bonded labour (Sahariya/Juang, Haliya); low literacy; land alienation (Baiga); no baseline surveys.
\item Scheme for Development of PVTGs (state plans, 100\% central funding); NAC 2013 (Sonia Gandhi) --- ST status to all PVTGs, rights-based approach, single-window, regular census/baseline surveys.
\end{itemize}
\end{QuickRevision}

\subsection{Geographical Distribution of Tribes}

\begin{DataFact}
\begin{itemize}
\item \textbf{2011 Census}: tribal population $\approx$ \textbf{10.42 crore}, i.e. a little over \textbf{8.5--8.6\%} of India's population.
\item Per the draft \textbf{National Tribal Policy}, there are $\approx$ \textbf{698} Scheduled Tribes ($\approx$ \textbf{705} per the Census).
\item \textbf{Odisha} has the largest number of distinct tribes --- \textbf{68}.
\item \textbf{Nine tribes} constitute about \textbf{60\%} of the total tribal population: \textbf{Bhil, Gond, Santhal, Oraon, Mina, Munda, Khond, Ho and Naga}. Eight of these are in central India (Chhattisgarh, Odisha, Jharkhand, Andhra); only the \textbf{Naga} are from the North-East.
\item Some tribes are found across five-six states; others are confined to a specific district of a single state.
\end{itemize}
\end{DataFact}

\begin{KeyConcept}
\begin{itemize}
\item \textbf{L. P. Vidyarthi} (the famous Indian anthropologist) divided India into \textbf{five zones} on the basis of tribal distribution:
\item 1. \textbf{Himalayan region} --- (a) \textbf{North-east Himalayan} (the seven sisters + Sikkim: Assam, Meghalaya, Arunachal, Manipur, Tripura, etc.) with tribes like \textbf{Naga, Garo, Khasi, Lepcha}; (b) \textbf{Central Himalayan} (the Tarai/Uttarakhand) with the \textbf{Khasa} and \textbf{Bhotiya}; (c) \textbf{North Himalayan} (Jammu \& Kashmir, Himachal) with the \textbf{Gaddi, Bakarwal, Gujjar}.
\item 2. \textbf{Middle India} --- Madhya Pradesh, Chhattisgarh, Odisha, Jharkhand and Bihar, where most of the tribal population (over 5--5.5 crore) lives.
\item 3. \textbf{Western India} --- Rajasthan, Gujarat, Maharashtra, with tribes like \textbf{Bhil, Mina, Maldhari}.
\item 4. \textbf{Southern India} --- Karnataka, Andhra, Tamil Nadu, Kerala, with tribes like \textbf{Chenchu, Toda, Kurumba}.
\item 5. \textbf{Island region} --- Andaman \& Nicobar and Lakshadweep: \textbf{Jarawa, Onge, Andamanese, Sentinelese} (Andaman) and the \textbf{Moplah}/Muslim tribes of Lakshadweep.
\end{itemize}
\end{KeyConcept}

\begin{ExampleBox}
\begin{itemize}
\item \textbf{Khasa} (Central Himalayan) --- a famous group practising \textbf{polyandry} (one woman, multiple husbands). They justify it by claiming descent from the \textbf{Pandavas} (who practised polyandry), but the real reason is the \textbf{scarcity of fertile land} in the hills: if three sons each marry separately and divide land, the holding shrinks each generation (1/3, then 1/9...), so polyandry keeps the land undivided and viable. \textbf{Polyandry is an adaptation to ecology, not merely a cultural practice.}
\item \textbf{Onge} (Andaman) --- studied by \textbf{Radcliffe-Brown} (of structural-functionalism).
\item \textbf{Muslim tribes} --- Lakshadweep's tribes are mainly Muslim (Moplah); in J\&K the \textbf{Bakarwal} and \textbf{Gujjar} are Muslim tribes. (The \textbf{Asifa} case victim belonged to the Bakarwal pastoral community.) Muslim STs get reservation --- unlike the separate question of Muslim/Christian SCs.
\end{itemize}
\end{ExampleBox}

\section{Tribes and the Colonial State}

\subsection{Three Phases of Tribe--State Interaction}

\begin{KeyConcept}
\begin{itemize}
\item The interaction of tribes with the state can be seen in \textbf{three phases}: \textbf{pre-colonial}, \textbf{colonial}, and \textbf{post-colonial} (post-independence).
\item \textbf{Pre-colonial period} --- interaction was very limited: tribes were given complete freedom and left \textbf{ungoverned}; early feudal lords paid little attention, recruiting them only for the military when needed. No taxation or laws were imposed on them.
\item Reason: tribal (forest) areas lay \textbf{between two kingdoms} and served as a \textbf{buffer zone} --- the kingdoms got free protection (any army passing through the forest first faced the tribals) and safe passage, so they maintained harmonious relations and never taxed them.
\end{itemize}
\end{KeyConcept}

\begin{KeyConcept}
\begin{itemize}
\item \textbf{Colonial period} --- the idea of the state changed radically: the British \textbf{unified India} under a single administration, which was \textbf{antithetical to tribal society} (lineages, clans, descent groups, kinship-based politics controlled by relatives).
\item It was a \textbf{despotic state} (the idea of \textbf{Karl Wittfogel}): a strong, centrally governed administrative machinery for \textbf{tax collection}, backed by a large military and police force.
\item For the \textbf{first time}, tribals tasted the \textbf{imposition of direct rule}.
\end{itemize}
\end{KeyConcept}

\subsection{The Colonial State: Three Aspects}

\begin{KeyConcept}
\begin{itemize}
\item British rule with respect to tribes can be summarised under \textbf{three aspects}:
\item 1. \textbf{Classification}.
\item 2. \textbf{Policy of segregation}.
\item 3. \textbf{Subjugation} --- which ultimately led to several \textbf{tribal revolts} (a cyclical relationship: these aspects triggered revolts, which led to more of the same).
\end{itemize}
\end{KeyConcept}

Policy-making is always dictated by the perceived interests of those in power. As long as tribes occupied remote regions and did not interfere with the rulers' autonomy, they were \textbf{left alone}. But when those inhospitable regions were noticed as \textbf{reservoirs of valued goods} (forest, minerals, mines --- the raw materials of British industry), the tribes began to be seen as \textbf{impediments to revenue generation}, and suppression, killing and genocide followed.

\subsection{Classification and the Criminal Tribes Act}

British administrators (governor-generals, commissioners) were highly educated (Oxford, Cambridge) and combined \textbf{subjugation} with \textbf{academic curiosity} --- colonial ethnographers like \textbf{Herbert Risley} and \textbf{J. H. Hutton} were themselves East India Company employees. This curiosity produced \textbf{classification} and several categories of tribe, most infamously the \textbf{'criminal tribes'} (now called \textbf{denotified tribes}).

\begin{ExampleBox}
\textbf{Hutton} entered the villages of the \textbf{Angami Naga} (infamous for \textbf{head-hunting} --- cutting and keeping an enemy's head as a trophy) to suppress the practice; even as he attacked them he collected their trophies and documented them --- suppression fused with academic study.
\end{ExampleBox}

\begin{GovtPolicy}
\begin{itemize}
\item \textbf{Criminal Tribes Act, 1871} --- its background lies in the \textbf{gypsies of Europe} and the \textbf{Industrial Revolution}: with the rise of a settled (sedentary) lifestyle as the norm, nomads were treated with \textbf{suspicion} (part of \textbf{Victorian ethics} that travelled to India --- though India never saw nomadism as a threat, e.g. Shankaracharya's travels).
\item First came the \textbf{Habitual Offenders Act, 1869}, giving authorities enormous power to regulate 'habitual offenders' (arrest, house arrest, punishment, questioning).
\item Combining this with the \textbf{caste system} and the idea of \textbf{race} (biological inheritance of traits --- just as some 'races' were deemed feeble-minded, criminality was treated as an inherited trait), the Criminal Tribes Act \textbf{branded whole communities criminal by birth}.
\item It \textbf{branded 150 communities criminal by birth}, made it compulsory for suspected tribes to be \textbf{registered}, put them under \textbf{strict surveillance}, curtailed their movements, and turned them into \textbf{scapegoats} --- a novel idea: a person labelled criminal from the day of birth irrespective of any act.
\end{itemize}
\end{GovtPolicy}

\begin{KeyConcept}
\textbf{Race is a 'folk taxonomy'} --- how a society categorises the outer world. In simple tribal society it is just \textbf{'us versus them'}; as society grows more complex, categories grow more complex (aided by evolutionism, which placed societies on a ladder of development --- some even called \textbf{'homo monstrosus'}, i.e. not fully human, monster-like). Racial classification was then arranged into a hierarchy and tasks allotted accordingly.
\end{KeyConcept}

\begin{CriticalPoint}
\begin{itemize}
\item Another reason for the Criminal Tribes Act was political: the \textbf{1857 First War of Independence} (which rural/tribal communities supported) and the \textbf{thagi (thuggee)} movement --- convoys passing through forests were robbed and British personnel killed. Communities that opposed British rule were therefore \textbf{labelled criminal} --- a tool to suppress revolt.
\item Almost \textbf{all} the communities branded criminal lived \textbf{near the main settlements/centres of British rule} (implemented first in Bengal, then extended to the North-West like Punjab --- where the revolts were).
\end{itemize}
\end{CriticalPoint}

\begin{GovtPolicy}
\begin{itemize}
\item \textbf{Post-independence}: on the recommendation of the \textbf{Iyengar Committee}, the Criminal Tribes Act was \textbf{repealed} by the \textbf{Criminal Tribes Laws Repeal Act, 1952}, replaced by the \textbf{Habitual Offenders Act, 1952} (blame shifted from the whole community to the individual, and only after repeat offences).
\item The \textbf{Dhebar Committee} suggested dropping the derogatory 'ex-criminal tribe' label in favour of \textbf{denotified tribes} (Hindi: \textbf{vimukta jatis}).
\item Yet the \textbf{stigma of criminality persists} --- these communities are still treated as criminal by local police/administration and remain marginalised.
\end{itemize}
\end{GovtPolicy}

\subsection{Policy of Segregation}

\begin{KeyConcept}
\begin{itemize}
\item The \textbf{policy of segregation and separation} was followed in parts of India with \textbf{no trading/commercial interest} --- mainly the \textbf{North-East and South India} --- where the British did not want to spend energy or money.
\item Although the British kept their hold on these far-off regions (appointing subordinate kings who treated them as paramount power), they did \textbf{not interfere} in tribal affairs.
\item Example: the \textbf{Naga Hills} had little intrusion until opened up during \textbf{World War II} (when Japan entered \textbf{Imphal and Kohima}, the North-East people supported the British and opened their village pathways to the army).
\item The British helped set up the \textbf{Naga Hill District Council} to deal with outsiders --- later renamed the \textbf{Naga National Council} under the leadership of \textbf{Angami Zapu Phizo} --- the seed of Nagaland's ethno-nationalism.
\item This produced the \textbf{Inner Line Permit (ILP)} --- a special permit (a 'visa') still required today to enter the North-East; \textbf{outsiders are not allowed to possess land/property} in the North-East (and in some other areas like East Godavari district).
\end{itemize}
\end{KeyConcept}

\begin{ComparisonBox}
\begin{itemize}
\item Under this segregation, the British divided tribal areas into \textbf{Excluded Areas} and \textbf{Partially Excluded Areas} (Government of India Act 1935) --- these later became the \textbf{Fifth Schedule} and \textbf{Sixth Schedule} areas.
\item This is a clear colonial carry-over: just as the 1935 Act's list of tribes became the basis of the Scheduled Tribes list, the Excluded/Partially Excluded areas became the Fifth/Sixth Schedule areas.
\end{itemize}
\end{ComparisonBox}

\begin{CriticalPoint}
\begin{itemize}
\item The segregation was largely \textbf{lip service} (especially in mainland India): land-transfer acts that should have protected tribal land failed because of three combined exploiters --- the \textbf{simplicity of tribes}, the \textbf{entry of money-lenders}, and \textbf{corrupt officers} --- which pushed tribes into marginalisation and poverty.
\item This is why tribal struggles are for \textbf{Jal, Jungle and Zameen} (water, forest and land) --- the three things snatched away from them.
\end{itemize}
\end{CriticalPoint}

\subsection{Subjugation and Tribal Revolts}

\begin{KeyConcept}
\begin{itemize}
\item \textbf{Subjugation} = active use of force, leading to bloody wars and brutal suppression of anyone challenging the British Raj --- through \textbf{Saam, Daam, Dand, Bhed} (all means) --- intensified by the British prejudice that tribals were \textbf{animal-like creatures} (violence in remote tribal areas went unnoticed).
\item Later, \textbf{Gandhi} mobilised people nationally, and the \textbf{Congress}, in a late session, passed a resolution that post-independence \textbf{tribal rights and culture would be preserved} and tribes given freedom to live as they liked.
\end{itemize}
\end{KeyConcept}

\begin{KeyConcept}
\begin{itemize}
\item \textbf{Major factors that led to tribal unrest:}
\item 1. \textbf{Infiltration of outsiders} (merchants, money-lenders, police).
\item 2. \textbf{Imposition of external rules} --- civil and criminal laws, and \textbf{monetization} of a non-monetized economy (money-lenders gave loans without security, then trapped the tribals into plantation labour to repay; taxes/lagaan were alien to tribal society).
\item 3. \textbf{Disproportionate impoverishment and marginalization} --- the colonial brunt was faced by all, but tribals faced the most severe attack (treated like animals); between \textbf{1891 and 1901} the tribal population \textbf{declined substantially}.
\end{itemize}
\end{KeyConcept}

\begin{ExampleBox}
\begin{itemize}
\item \textbf{Santhal uprising (1855)} --- against merchants, money-lenders and police imposing their rule; approximately \textbf{30,000 Santhals} were killed.
\item \textbf{Mina uprising (Rajasthan)} under the leadership of \textbf{Govind Guru} --- against the \textbf{begari/begar} system (bonded labour: a loan of \rupee 500 binding one to a lifetime of work).
\end{itemize}
\end{ExampleBox}

\subsection{Impact of Colonial Rule on Tribal Life}

\begin{KeyConcept}
\begin{itemize}
\item 1. \textbf{Introduction of private property} --- broke down the \textbf{communal mode of production} and communal ownership of resources (forest owned collectively, not by an individual/family), eroding the collective solidarity.
\item 2. \textbf{Land alienation and 'tenantisation'} --- tribes became \textbf{tenants in their own land}, leading to \textbf{peasantization} (hunter-gatherers forced into cultivation) and \textbf{depeasantization} (cultivators reduced to daily-wage/contract labour) --- through large-scale \textbf{land transfer from adivasis to non-adivasis}.
\item 3. \textbf{Commercialization of forest} --- further marginalisation. Verrier Elwin's \textbf{Nature--Man--Spirit (NMS) complex}: for the tribe the forest is home and the dwelling of their god, and their economy, social life and culture are forest-centric --- so removing them is \textbf{ethnocide} (killing the culture), not just economic loss.
\item Example: \textbf{Niyamgiri} --- the \textbf{Dongria Kondh} tribe opposed the allocation of the Niyamgiri forest to \textbf{Vedanta} (and won the case); they believe the trees are the home of their god \textbf{Niyam Raja} and cutting them would bring a divine curse.
\item 4. \textbf{Monetization} --- money replaced \textbf{barter} and introduced an \textbf{impersonal medium of exchange} (in barter, exchange is embedded in social obligation --- you cannot exchange with anyone; in money exchange, identity does not matter). This meant the \textbf{general disappearance of non-market mechanisms of social control of resources and allocation of goods}.
\item 5. \textbf{Entry of class} --- monetization created poverty, confusion and \textbf{class differentiation} within tribes (some with more harmonious contact with the mainstream became elites).
\end{itemize}
\end{KeyConcept}

\subsection{Soft Power and the Overall Impact}

\begin{KeyConcept}
\begin{itemize}
\item Alongside classification, segregation and subjugation, the colonial power also used \textbf{soft power}: it \textbf{funded and sent Christian missionaries} to the tribal areas --- to 'civilise' criminal tribes and to spread Christianity (most successfully in the North-East and Eastern India) --- while simultaneously using brutal violence.
\item \textbf{Overall impact}: the transformation of \textbf{Adivasi dignity} and a \textbf{threat to tribal identity}. Tribes became \textbf{'the other'} --- whose economic non-conformity came in the way of the exploitation of natural resources.
\end{itemize}
\end{KeyConcept}

\begin{CriticalPoint}
\begin{itemize}
\item This threat to tribal identity did \textbf{not end at independence} --- poverty, marginalisation and impoverishment consolidated the ethnic identity and produced \textbf{two responses}:
\item 1. \textbf{Ethno-nationalism} --- tribal parties and movements asserting identity (e.g. \textbf{Jharkhand Mukti Morcha}).
\item 2. \textbf{Naxalism (left-wing extremism)} --- when the state/rulers themselves become exploiters, such ideologies find fertile ground. Today Naxalism is concentrated in the \textbf{tribal belts} --- the \textbf{'red corridor'} of Odisha, Andhra and West Bengal.
\end{itemize}
\end{CriticalPoint}

\begin{QuickRevision}
\begin{itemize}
\item Three phases: pre-colonial (buffer zones, no taxation, military recruitment) $\rightarrow$ colonial (unified despotic state, direct rule) $\rightarrow$ post-colonial.
\item Three aspects of British rule: classification, segregation, subjugation $\rightarrow$ tribal revolts.
\item Criminal Tribes Act 1871 (from Habitual Offenders Act 1869 + caste + race) --- 150 communities branded criminal by birth; repealed 1952 (Iyengar Committee) $\rightarrow$ Habitual Offenders Act 1952; renamed 'denotified tribes' (vimukta jatis) by Dhebar Committee; stigma persists.
\item Segregation: NE/South India lip service; Excluded/Partially Excluded areas $\rightarrow$ 5th/6th Schedule; Inner Line Permit; Naga National Council (Zapu Phizo).
\item Subjugation: Saam-Daam-Dand-Bhed; factors --- outsiders, external rules/monetization, disproportionate impoverishment; Santhal uprising 1855 (\textasciitilde{}30,000 killed), Mina uprising (Govind Guru, begari).
\item Impacts: private property $\rightarrow$ breakdown of communal mode; land alienation $\rightarrow$ peasantization/depeasantization; forest commercialization $\rightarrow$ ethnocide (Niyamgiri/Dongria Kondh); monetization $\rightarrow$ impersonal exchange; class penetration; Christian missionaries as soft power.
\item Post-colonial: threat persists $\rightarrow$ ethno-nationalism (Jharkhand Mukti Morcha) or Naxalism (red corridor).
\end{itemize}
\end{QuickRevision}

\section{Post-Colonial Administration --- Isolation, Assimilation and Integration}

\subsection{The Post-Independence Dilemma}

\begin{KeyConcept}
\begin{itemize}
\item Post-independence, the Indian government faced a dilemma over how to manage the tribes. Two major approaches contended:
\item 1. The \textbf{isolationist approach} --- advocated by \textbf{Verrier Elwin}.
\item 2. The \textbf{assimilationist approach} --- advocated by \textbf{G. S. Ghurye (Govind Sadashiv Rao Ghurye)}.
\item Both, in their own way, were really aiming at the same end --- ultimately making the tribe part of Hindu society --- but by different routes.
\end{itemize}
\end{KeyConcept}

\subsection{Verrier Elwin and the Theory of Isolation}

\textbf{Verrier Elwin} was a missionary who came to India, worked in tribal areas, came in contact with Mahatma Gandhi, and assisted the freedom struggle; post-independence he was appointed \textbf{tribal advisor to Jawaharlal Nehru}. His core finding: whenever tribes come into contact with the outside world they are \textbf{exploited and left marginalised} --- hence his theory of isolation.

\begin{KeyConcept}
\begin{itemize}
\item Elwin advocated establishing \textbf{national-park-like structures}, isolated from the outside world, where entry is not possible without permission.
\item He wrote the book \textbf{'The Baiga'} (on the Baiga tribe), arguing the government should adopt this national-park approach.
\item Reason: on contact with the mainstream, tribes are exploited by \textbf{money-lenders, zamindars and corrupt police} --- so isolation lets them live happily.
\item Within this 'national park': allow the \textbf{traditional tribal councils (panchayats)} and recognise the authority of the \textbf{village headman/chief}; any non-tribal settlement (hospital, school, bank) requires a \textbf{government license and the tribal head's permission}; and \textbf{no missionary of any religion} should be allowed in, so tribal life is not broken up.
\item In short: minimise the tribe's contact with outsiders.
\end{itemize}
\end{KeyConcept}

\subsection{Elwin's Four Classes of Tribes ('The Aboriginals')}

In \textbf{'The Aboriginals' (1943)}, Elwin divided tribes into \textbf{four classes} on the basis of their cultural development / amount of cultural contact with outsiders.

\begin{KeyConcept}
\begin{itemize}
\item \textbf{Class 1} --- the purest of pure tribal groups: completely isolated from the mainstream, about \textbf{2--3 million} persons, with their \textbf{religion alive}, \textbf{tribal organisation unimpaired} and a healthy organisation of tribal life.
\item \textbf{Class 2} --- tribes that have come into contact with the mainstream Hindu society but retain much of their tribal characteristics. Their life has become more \textbf{individualistic} and less honest and less simple (communal sharing/exchange is less prevalent).
\item \textbf{Class 3} --- the \textbf{largest section} (about \textbf{four-fifths} of the tribal population): \textbf{tribal in name but Hindu in practice}, absorbed into the \textbf{lower rung} of the Hindu hierarchy (informally treated as lower castes). Causes: British exploitation (uprooted, migrated to cities, unskilled manual labour, poverty, illiteracy, the vicious cycle of poverty and indebtedness) and conversion/absorption by \textbf{Christian and Hindu missionaries}.
\item \textbf{Class 4} --- those who have \textbf{'won the battle of culture'}: the old \textbf{tribal aristocracy} --- \textbf{Raj Gonds (Gond Rajas), great Bhils, Naga chieftains, wealthy Santhals, highly intellectual Mundas} --- economically well-off and properly recognised by the mainstream; they retain their old tribal names and tribal religion despite adopting the Hindu faith and modern lifestyle.
\end{itemize}
\end{KeyConcept}

\begin{KeyConcept}
\begin{itemize}
\item Elwin's stated aim: \textbf{the whole 'aboriginal problem' is to enable tribesmen of the first and second classes to advance directly to the fourth class, without suffering the despair and degradation of the third}.
\item Hence he advocates \textbf{isolation, especially during the transition period} --- first make them economically powerful, and expose them to competition only when they are ready.
\item This is really a gentler path to the same assimilation --- but avoiding the painful third stage.
\end{itemize}
\end{KeyConcept}

\subsection{Ghurye and the Theory of Assimilation}

\textbf{G. S. Ghurye} argued for assimilation in his book \textbf{'The Scheduled Tribes'}. He was backed by right-wing/Indological scholars, notably \textbf{A. V. Thakkar Bappa} (head of a Hindu mission --- like the Christian missionaries --- running schools such as Saraswati Shishu Mandir), who \textbf{criticised Elwin's isolation as a 'colonial hangover'} / anti-national (adopting the British policy of segregation).

\begin{KeyConcept}
\begin{itemize}
\item Like Elwin, Ghurye divides the tribes into \textbf{three sections}:
\item \textbf{Section 1} --- the \textbf{Raj Gonds} and similar tribes who have \textbf{successfully fought the battle of cultural contact} and are recognised as members of fairly high status (even proclaiming \textbf{Kshatriya} status) in Hindu society.
\item \textbf{Section 2} --- the \textbf{large mass of tribes} that has been \textbf{partially Hinduized} through closer contact with Hindus (plain peasants/villagers/rural society).
\item \textbf{Section 3} --- the \textbf{hill sections} that exhibit great resistance to the alien culture --- completely isolated, treating Hindu culture as alien, shy of contact and maintaining distance.
\item Most tribes are in this third category.
\end{itemize}
\end{KeyConcept}

\begin{KeyConcept}
\begin{itemize}
\item Ghurye's conclusion: \textbf{'tribal people in India are backward Hindus, differing only in degree.'}
\item Their \textbf{backwardness is due to imperfect/partial assimilation} into Hindu society --- had they been fully Hinduized they would not be backward, marginalised or poor.
\item Partial Hinduization creates a conflict: they are \textbf{bilingual} (dominant language outside, own language at home) --- needing the mass language for survival while protecting identity through their own language.
\item Therefore the remedy is \textbf{complete assimilation} --- create paths so their assimilation becomes more organised and sorted.
\end{itemize}
\end{KeyConcept}

\subsection{Nehru's Policy of Integration}

\begin{KeyConcept}
\begin{itemize}
\item The government chose \textbf{neither assimilation nor isolation}, but the \textbf{policy of integration} --- the stand of \textbf{Jawaharlal Nehru} and his government.
\item In fact, after fierce criticism, \textbf{Elwin himself changed his stance} and proposed the integration policy (though it is normally attributed to Nehru).
\item Nehru's reasoning: we should \textbf{not isolate} them (treating tribes like \textbf{museum artifacts} --- static specimens that anthropologists visit and write about), and we should \textbf{not assimilate} them (destroying their prestigious, virtuous culture).
\item Nehru's famous line: \textbf{'We should not try to make them a second copy of ourselves'} --- instead, provide the conditions for them to live their life to the fullest, retain their culture, and enjoy it, while giving them basic necessities (health, education, electricity) without hampering their culture, and \textbf{empower} them to become self-dependent.
\item Criticism: this was \textbf{lip service} --- in letter but not in spirit --- because most colonial policies were carried forward as they were, and tribes continued to be exploited.
\end{itemize}
\end{KeyConcept}

\subsection{The Tribal Panchsheel (Five Principles)}

\begin{GovtPolicy}
\begin{itemize}
\item Nehru's policy of integration is expressed in the \textbf{Tribal Panchsheel} --- \textbf{five fundamental principles for tribal development}, aimed at protecting indigenous/tribal culture:
\item 1. \textbf{People should develop along the line of their own genius} --- not imposing anything on them; encouraging their own traditional art and architecture and providing a market for their products (e.g. if they are good at weaving, let them develop that).
\item 2. \textbf{Tribal rights on lands and forests should be protected.}
\item 3. \textbf{Train and build up a team of their own people} to do the work of administration and development --- avoiding too many outsiders in their territory; the ultimate aim is an army of workers from among them (few outsiders want to serve in tribal areas, e.g. the film \emph{Newton}).
\item 4. \textbf{Do not over-administer those areas or overwhelm them with a multiplicity of schemes} --- rather, work through their own social and cultural institutions (their traditional ways of managing resources and resolving disputes).
\item 5. \textbf{Judge results not by the statistics of money spent, but by the quality of human character evolved} --- measure impact, not expenditure.
\end{itemize}
\end{GovtPolicy}

\begin{CriticalPoint}
The third principle was never realised because of \textbf{tribal/dalit elitism}: those who obtain reservation benefits often associate themselves with the higher group and refrain from working for the empowerment of their own community --- the same problem that weakens reservation benefits for Dalits as a whole.
\end{CriticalPoint}

\subsection{Schemes of Integration (TSP, Fifth/Sixth Schedule, PESA)}

\begin{GovtPolicy}
\begin{itemize}
\item Under the policy of integration, the government designed several schemes/structures:
\item \textbf{Tribal Sub-Plan (TSP)} --- a directly asked exam question (2017--18, 10-mark). It was designed in the \textbf{Fifth Five-Year Plan} on the recommendation of the anthropologist \textbf{S. C. Dube} (with help from \textbf{L. P. Vidyarthi}).
\item \textbf{Fifth Schedule areas} and \textbf{Sixth Schedule areas}.
\item \textbf{PESA --- Panchayats (Extension to Scheduled Areas) Act}, which falls under the Fifth Schedule areas.
\item (These are covered in polity/Laxmikanth.)
\end{itemize}
\end{GovtPolicy}

\begin{QuickRevision}
\begin{itemize}
\item Post-independence dilemma: isolation (Elwin) vs assimilation (Ghurye).
\item Elwin: national-park model (The Baiga); no outside entry; tribal councils + headman authority; no missionaries. The Aboriginals (1943) --- 4 classes (Class 1 purest; Class 2 contact but tribal; Class 3 tribal-by-name absorbed in lower rung, \textasciitilde{}4/5; Class 4 aristocracy). Aim: Class 1--2 $\rightarrow$ Class 4 directly, skipping Class 3.
\item Ghurye (The Scheduled Tribes): 3 sections (Raj Gonds high-status; partially Hinduized mass; resistant hill sections). 'Tribes are backward Hindus'; backwardness = imperfect assimilation $\rightarrow$ complete assimilation.
\item Nehru: integration --- not museum, not second copy of ourselves; but lip service in practice.
\item Tribal Panchsheel (5 principles): own genius; protect land/forest rights; build own personnel; don't over-administer (work through their institutions); judge by human character, not money spent.
\item Schemes: Tribal Sub-Plan (5th FYP, S. C. Dube), Fifth/Sixth Schedule, PESA.
\end{itemize}
\end{QuickRevision}

\section{Ethnicity, Ethnic Groups and Ethno-Nationalism}

\subsection{Nation, State and Ethnicity: The Basic Concepts}

The term \textbf{ethnicity} is derived from \textbf{ethnos}, meaning \textbf{nation}. Understanding ethnicity therefore begins with the difference between a nation and a state.

\begin{KeyConcept}
\begin{itemize}
\item \textbf{Nation} --- a group of people with \textbf{shared culture and belief}, who use this shared culture for \textbf{political mobilisation} (i.e. there is political awareness/consciousness). Shared culture alone does not make a nation; the political identity/consciousness is essential.
\item \textbf{State} --- a \textbf{political/geopolitical entity} with a \textbf{territory} of its own.
\item A \textbf{nation may or may not have a territory}; when a politically aware group with shared culture gets its own territory it becomes a \textbf{nation-state} (e.g. \textbf{Japan}).
\item Examples of nations without states: the \textbf{Kurds} (fighting for a territory), and \textbf{Khalistan} (a nation with cultural belief and political consciousness used to mobilise, demanding a state).
\item A state can be multicultural, pluralistic or monocultural; when the majority of a state's people share one culture it becomes a nation-state.
\end{itemize}
\end{KeyConcept}

\begin{KeyConcept}
\begin{itemize}
\item \textbf{Ethnicity} (definition by \textbf{Thomas Hylland Eriksen}): \textbf{the relationship between groups of people whose members consider themselves distinct, and who are often related or ranked hierarchically within a society.}
\item Two essential elements: (1) an \textbf{awareness of distinctiveness}, and (2) the groups are \textbf{ranked hierarchically} (this ranking is what drives ethnic interaction, as groups try to protect or improve their rank).
\end{itemize}
\end{KeyConcept}

\subsection{Defining Ethnicity and the Ethnic Group}

\begin{KeyConcept}
\begin{itemize}
\item \textbf{Collective identity} is based on some common characteristics --- which can be \textbf{objective} (region, language, religion, caste --- clear, discrete categories based on shared experience/history) or \textbf{subjective} (by self-ascription --- a sense of belonging even without actually belonging).
\item \textbf{Ethnic group} = a \textbf{historical entity} whose members, in large part, \textbf{consider themselves alike} by virtue of certain \textbf{stable features} (language, race, culture, territory, etc.).
\item An \textbf{ethnic group can never exist in isolation} --- ethnic identity surfaces only when there is \textbf{contact} with another group and a \textbf{ranking/deprivation} against that group.
\item Example: Tamil speakers live happily until the state tries to impose Hindi as a national language --- the ranking/threat makes their linguistic ethnic identity surface with solidarity.
\end{itemize}
\end{KeyConcept}

\begin{ExampleBox}
\begin{itemize}
\item \textbf{Subjective/self-ascribed identity} --- \textbf{Jharkhand Mukti Morcha}: it aimed to liberate Jharkhand from Bihar, arguing tribal rights were not recognised; but not only tribals supported it --- several caste groups supported it for their \textbf{vested interest} (in a new state they would become more powerful). This is called \textbf{'pseudo-tribalism'} or \textbf{'tribalism without tribes'} --- a \textbf{contextual/situational identity} assumed, not objective.
\item Likewise, leaders of \textbf{Dalit assertion} and \textbf{feminist} movements are often not themselves from the marginalised group (e.g. \textbf{Gandhi} fighting untouchability though not a Dalit) --- they \textbf{self-ascribe} to the movement's identity.
\item A \textbf{situational} example: abroad, a Bihari and a person from Madhya Pradesh may bond over a shared language even though their identities differ at home.
\end{itemize}
\end{ExampleBox}

\begin{KeyConcept}
\begin{itemize}
\item According to \textbf{Eriksen}: ethnic groups can never exist in isolation; very often there is a \textbf{hierarchy within society}, with some groups enjoying \textbf{more status and greater material rewards} than others.
\item For ethnic groups to develop there \textbf{must be some contact between groups} that entertains the idea of each other as being culturally distinct.
\item When contact is \textbf{not beneficial} (deprivation, exploitation), it produces ethnic identity; when this identity is used for \textbf{political mobilisation}, it becomes \textbf{ethno-nationalism} (a sense of national identity). The politically aware group may then take democratic routes (protest, discussion) or turn to violence (civil war, secessionist movements, Naxalism).
\end{itemize}
\end{KeyConcept}

\subsection{Criteria of an Ethnic Group (M. E. Brown and Weber)}

\begin{KeyConcept}
\begin{itemize}
\item \textbf{Michael E. Brown} gave \textbf{six criteria} for a group to be an ethnic group:
\item 1. \textbf{Name} --- a group must have a name to identify it.
\item 2. \textbf{Common ancestry} --- real, assumed or putative (e.g. the \textbf{Yadavs} believe they descend from Krishna / 'Yaduvanshi'); the belief itself creates solidarity.
\item 3. \textbf{Shared belief system} --- at least on contextual issues (e.g. in the \textbf{farmers' protest}, people from very different backgrounds --- Bollywood/Punjabi stars, capitalists, money-lenders, intermediaries, genuine farmers --- share a belief in that context despite their other conflicts).
\item 4. \textbf{Some degree of shared culture} --- shared belief without shared culture is not enough (I may believe Blacks in the US are wrongly exploited and sympathise, but I will never be part of that ethnic group because I do not share the culture and would not be accepted as a member).
\item 5. \textbf{Sense of attachment to a particular territory} --- common but not fully accepted, because some ethnic groups have no territorial attachment (e.g. \textbf{Dalit emancipation}).
\item 6. \textbf{Self-awareness} --- most members must believe they constitute an ethnic group (awareness 'in mind, heart and soul' that I am a member of this group in a particular role).
\end{itemize}
\end{KeyConcept}

\begin{KeyConcept}
\textbf{Max Weber's} similar criteria for an ethnic group: (1) \textbf{shared ancestry} (real or assumed), and (2) \textbf{shared experience} --- a shared history (language, historical past, colonisation, deprivation, shared territory). This common belief creates strong bonds between members.
\end{KeyConcept}

\subsection{Types of Ethnicity (Eriksen)}

\begin{KeyConcept}
\begin{itemize}
\item \textbf{Eriksen's classification} of types of ethnic interaction (US-based but valid for India):
\item 1. \textbf{Modern migrants} --- recent immigrants and their descendants (e.g. in the US); they suffer racism, are aware of being weaker than the dominant white/Caucasoid group, but are \textbf{unlikely to demand their own state} (little political mobilisation).
\item 2. \textbf{Indigenous population} --- original inhabitants living there before colonisation; \textbf{politically powerless} (relative isolation, lack of education/financial empowerment, under-represented in national government), and \textbf{partly integrated into the dominant society at the lower hierarchy} (like Elwin's 'Class 3'). E.g. the weaker tribes --- \textbf{Maldhari} (Gujarat), \textbf{Bonda}, \textbf{Toda}, the \textbf{PVTGs} --- exploited but too few/powerless to resist or demand a state.
\item 3. \textbf{Proto-nations / ethno-nationalists} --- ethnic groups \textbf{actively seeking an independent state} --- e.g. the \textbf{Nagas} (Naga National Council / NSCN), \textbf{Meiteis} (before merger), several North-East secessionist groups (supported by China/Myanmar), the \textbf{Kurds}, \textbf{Tamils in Sri Lanka (LTTE)}, and the \textbf{Balochis}.
\item 4. \textbf{Post-slavery minorities} --- descendants of slaves (e.g. \textbf{Blacks/Afro-Americans} in the US, the \textbf{Siddhi} community in Gujarat, people of African origin in Malaysia); now citizens but still suffering racism/degradation $\rightarrow$ \textbf{identity politics}.
\end{itemize}
\end{KeyConcept}

\subsection{Multiculturalism vs Pluralism}

\begin{ComparisonBox}
\begin{itemize}
\item \textbf{Multiculturalism} --- many cultures living in the same society, held together by \textbf{tolerance} (stereotypes persist and prevent interaction; there is merely law and order).
\item \textbf{Pluralism} --- different cultural groups living together with \textbf{constructive dialogue}: they sit together, discuss each religion's problems, modify and appreciate each other's traditions.
\item The aim is \textbf{pluralism} (the utopian end); in reality India is \textbf{multicultural} --- tolerance rather than constructive dialogue.
\item Example: in a pluralist society, Hindus and Muslims would sit and discuss their religions; in a multicultural society there is only tolerance backed by law and order.
\end{itemize}
\end{ComparisonBox}

\subsection{Hot and Cold Ethnicity (Fenton)}

\begin{KeyConcept}
\begin{itemize}
\item \textbf{Steve Fenton} described two types of ethnicity based on \textbf{degree of intensity}:
\item \textbf{Hot ethnicity} --- appeals to \textbf{blood and passion}, with \textbf{strong group loyalty}; mobilised in support of nationalist movements or where there is intense inter-group conflict (e.g. communal riots --- people take up swords, petrol bombs).
\item \textbf{Cold ethnicity} --- much \textbf{less impassioned}, involving \textbf{calculation and instrumentality} (utilitarian): the appeal is less emotional and hopes to derive \textbf{benefits} from the group. Examples: election-time caste mobilisation (building a statue/temple for a caste's saint, declaring a festival), and companies using emotional ethnic appeals to sell products (Amazon's Rakhi ad, MakeMyTrip's Diwali ad).
\item In cold ethnicity one supports a group \textbf{in hope of deriving benefits} (OBC/SBC promises in manifestos), not out of blood-and-passion loyalty.
\end{itemize}
\end{KeyConcept}

\begin{KeyConcept}
\textbf{Fenton's key point}: \textbf{ethnicity is never the only identity people have --- no one is full-time ethnic} (at least not in the same way in all settings). A person has \textbf{multiple ethnic identities} (Hindu, Bihari, Rajput) that surface in \textbf{different circumstances} --- which identity is expressed depends on context.
\end{KeyConcept}

\subsection{Approaches to Ethnicity: Primordial, Instrumentalist, Materialist}

\begin{KeyConcept}
\begin{itemize}
\item \textbf{Primordial approach} (oldest, dominant before the 1970s; \textbf{Clifford Geertz}): ethnic identity is \textbf{innate, fixed and permanent} --- each individual is \textbf{born into} an ethnic group (a culturally defined unit), and ethnicity serves a \textbf{fundamental human need for belonging and meaning}. One is aware of one's cultural distinctiveness from birth.
\item It invokes \textbf{'ancient hatred'}: fundamental cultural differences and divergent values between ethnic groups inevitably result in a \textbf{clash of cultures}, caused by \textbf{past animosity}.
\item \textbf{Flaw}: it cannot explain violence where there was no past animosity --- e.g. the \textbf{sons-of-the-soil movement} in Maharashtra (initially Marathis vs 'Madrasis'/South Indians, later against \textbf{Biharis} --- with whom there was no historical conflict).
\end{itemize}
\end{KeyConcept}

\begin{KeyConcept}
\begin{itemize}
\item \textbf{Instrumentalist approach} (\textbf{Fredrik Barth} --- 'ethnic boundaries'): ethnic identity is \textbf{instrumental/utilitarian} --- it expresses only when there is some utility.
\item \textbf{Ethnic boundaries are social rather than territorial}: individuals \textbf{selectively emphasise} those forms of cultural difference that are important to them, and the cultural features drawn upon are \textbf{not fixed but situationally defined}.
\item Which identity surfaces depends on context: in a Hindu--Muslim conflict the \textbf{Hindu} identity comes up (forgetting caste); in an intercaste conflict the \textbf{Rajput/Yadav} identity comes up; when a Bihari goes to Maharashtra, the \textbf{Bihari} identity unites all Bihari Hindus against Marathis.
\item This is an \textbf{oscillating equilibrium} --- a person is a member of multiple ethnic groups and expresses different identities in different situations. \textbf{Ethnic identities are the product of ascription and self-ascription.}
\end{itemize}
\end{KeyConcept}

\begin{KeyConcept}
\textbf{Materialist (Marxist) approach}: ethnicity emerges from \textbf{class relations}. Either (1) ethnic violence is due to \textbf{economic inequality} between groups (the deprived group revolts against the oppressor), or (2) ethnicity is a \textbf{tool used by the bourgeoisie/capitalists to divide the proletariat} into ethnic groups so they never unite and never make a revolution.
\end{KeyConcept}

\subsection{Ethnic Conflict and Collective Violence}

\begin{KeyConcept}
\begin{itemize}
\item \textbf{Conflict $\neq$ violence.} Conflict can be expressed democratically through the avenues provided by democracy (discussion, dissent). Violence erupts when conflict is \textbf{suppressed by force} by an authoritarian/coercive state.
\item In an \textbf{ethnically plural society with freedom of expression}, some conflict on identity basis is \textbf{typically to be expected} (it is part of democracy --- active discussion). Issues discussed: which language should be used for education/employment; whether some groups should get rights; whether religious dress is allowed; affirmative action/reservation.
\item In an \textbf{authoritarian society} that keeps ethnic discontent under the surface, there is \textbf{sudden eruption of violence} (e.g. \textbf{Shaheen Bagh} --- violence erupted when right-wing groups entered and demanded clearance).
\end{itemize}
\end{KeyConcept}

\begin{KeyConcept}
\begin{itemize}
\item \textbf{Collective violence} (not individual) can take several forms:
\item 1. \textbf{Riots} --- between two ethnic groups on an \textbf{almost equal pedestal} (neither weak enough to be massacred), with state neutrality in doubt.
\item 2. \textbf{Pogrom / genocide} --- a \textbf{state-approved attack on helpless minorities} (e.g. the \textbf{Jews/Holocaust}, the \textbf{1984 anti-Sikh riots}).
\item 3. \textbf{Civil war} --- the \textbf{state itself becomes a combatant} (e.g. \textbf{Kurds} vs Turkey, \textbf{Tamils/LTTE} vs the Sri Lankan state, \textbf{Naga} secessionism).
\item 4. \textbf{Lynching} --- many kill one/few (mob justice; e.g. \textbf{Nido Tania}, an Arunachal Pradesh boy beaten to death by shopkeepers in Lajpat Nagar, 2015).
\item 5. \textbf{Terrorism} --- a few kill many (e.g. 'mujahideen' who know they will not return but kill many).
\end{itemize}
\end{KeyConcept}

\subsection{Causes of Ethno-Nationalism among Tribes}

\begin{KeyConcept}
\begin{itemize}
\item The reasons for the \textbf{rise of ethno-nationalism among tribes} (a directly asked exam question):
\item 1. \textbf{We-feeling and ethnocentrism} --- ethnic identity produces an 'in-group' feeling, leading to \textbf{ethnocentrism} and a \textbf{demand for more power}, mainly due to a \textbf{feeling of marginalisation} (they demand autonomy, forest rights, control over industry).
\item 2. \textbf{Historic factors} --- division of society on the basis of race, religion, language, region; \textbf{past animosity} (e.g. the \textbf{Bru/Reang} and \textbf{Meitei} cases, and the recent \textbf{Assam--Mizoram} border clash).
\item 3. \textbf{Political factors} --- lack of representation, no reservation/government jobs, feeling powerless; and \textbf{exclusionary nationalist ideologies} (e.g. the \textbf{Bodos} in Assam over citizenship; the \textbf{Darjeeling/Gorkhaland} agitation after Bengali was imposed as official language over Nepali speakers; also \textbf{Kodoland} and \textbf{Nagalim}).
\item 4. \textbf{Economic reasons} --- feeling of deficit/deprivation/backwardness, developmental deficit, resources snatched away and livelihood lost: \textbf{Jharkhand} (tribal resources feeding Bihar $\rightarrow$ Jharkhand Mukti Morcha), \textbf{Telangana} (Andhra spending concentrated in Rayalaseema).
\item 5. \textbf{Cultural factors} --- threat to identity and autonomy, plus \textbf{stereotypes/discrimination} and \textbf{language-related conflicts}.
\end{itemize}
\end{KeyConcept}

\begin{QuickRevision}
\begin{itemize}
\item Ethnicity (from 'ethnos' = nation) = relationship between groups that consider themselves distinct and are ranked hierarchically (Eriksen).
\item Ethnic identity surfaces only with contact + ranking/deprivation (not in isolation); objective vs subjective (self-ascribed) identity --- Jharkhand Mukti Morcha = 'pseudo-tribalism'.
\item Ethnic group criteria (M. E. Brown): name, common ancestry, shared belief, shared culture, territorial attachment, self-awareness; Weber: shared ancestry + shared experience.
\item Types (Eriksen): modern migrants, indigenous population, proto-nations/ethno-nationalists (Nagas, Kurds, LTTE, Balochis), post-slavery minorities (Blacks, Siddhi).
\item Multiculturalism (tolerance) vs pluralism (constructive dialogue); India aims for pluralism, is multicultural.
\item Hot vs cold ethnicity (Fenton); 'no one is full-time ethnic' --- multiple identities surface by context.
\item Approaches: primordial (Geertz --- innate/fixed, 'ancient hatred'), instrumentalist (Barth --- ethnic boundaries, situational), materialist (Marx --- class/divide-and-rule).
\item Collective violence: riots, pogrom/genocide, civil war, lynching, terrorism.
\item Causes of tribal ethno-nationalism: we-feeling/ethnocentrism; historic factors; political (representation, exclusionary ideology); economic (Jharkhand, Telangana); cultural (identity threat, language).
\end{itemize}
\end{QuickRevision}

